% Generated draft: requires mathematical and source review.
\documentclass[11pt]{article}
\usepackage[margin=25mm]{geometry}
\usepackage{fontspec}
\setmainfont{Latin Modern Roman}
\usepackage{amsmath,amssymb,mathtools,unicode-math}
\setmathfont{Latin Modern Math}
\usepackage{enumitem,needspace,xurl,hyperref}
\setlength{\parindent}{0pt}
\setlength{\parskip}{5pt}
\setlength{\emergencystretch}{3em}
\newcommand{\N}{\mathbb{N}}
\newcommand{\Z}{\mathbb{Z}}
\newcommand{\Q}{\mathbb{Q}}
\newcommand{\R}{\mathbb{R}}
\newcommand{\C}{\mathbb{C}}
\newcommand{\F}{\mathbb{F}}
\newcommand{\A}{\mathbb{A}}
\newcommand{\PP}{\mathbb{P}}
\newcommand{\E}{\mathbb{E}}
\newcommand{\eps}{\varepsilon}
\newcommand{\dd}{\,\mathrm{d}}
\begin{document}
\section*{UnsolvedMath ID: 30001717}
\textbf{Source problem code:} catalogue entry
\subsection{Definitions and mathematical statement}
Let \(q\) be a prime power and let \(G_q=\operatorname{GL}_n(\F_q)\). Interpret the irreducible characters \(X_1,\ldots,X_k\) as ordinary complex characters. Their tensor product has character \(g\mapsto\prod_{i=1}^k X_i(g)\). Write \(1_{G_q}\) for the trivial character and define its multiplicity by
\[
M_q(X_1,\ldots,X_k):=\left\langle X_1\otimes\cdots\otimes X_k,1_{G_q}\right\rangle_{G_q}
=\frac{1}{|G_q|}\sum_{g\in G_q}\prod_{i=1}^k X_i(g).
\]
Thus \(M_q\) is a nonnegative integer for each individual tuple.

The supplied conjecture concerns tuples described as \emph{generic}, characters described as \emph{split}, and an associated quiver dimension vector \(v\). The record does not define these qualifications or give the construction of the quiver and \(v\). Consequently, the following is a source-faithful schematic statement requiring those data before it becomes a fully specified conjecture.

For a family of generic tuples across the admissible prime powers \(q\), with the relevant character data and associated quiver dimension vector held fixed in a manner that remains unspecified, the conjecture asks whether there exists a polynomial \(P(t)\in\Z[t]\), independent of \(q\), such that
\[
M_q(X_1,\ldots,X_k)=P(q).
\]
It further asks whether \(P\) has positive coefficients when the characters are split, whether its leading coefficient is \(1\), and whether
\[
M_q(X_1,\ldots,X_k)\ne0
\quad\Longleftrightarrow\quad
v\text{ is a root of the associated Kac--Moody root system},
\]
\[
M_q(X_1,\ldots,X_k)=1
\quad\Longleftrightarrow\quad
v\text{ is a real root of that root system}.
\]
Here a real root means a root in the Weyl-group orbit of a simple root; the particular root system is not specified in the record.

\textbf{Unresolved qualifications.} The meaning of positive coefficients is not specified: it could mean nonnegative coefficients or strict positivity of some specified coefficients. The leading-coefficient assertion cannot apply to the zero polynomial, whereas the root criterion allows vanishing; the record does not explicitly restrict monicity to the nonzero case. It also does not specify whether the assertion of equality to \(1\) is pointwise in \(q\) or an assertion that \(P(t)\equiv1\). No resolution of these issues is imposed here.
\subsection{Short English statement}
The conjecture relates the multiplicity of the trivial character in a generic tensor product of irreducible characters of \(\operatorname{GL}_n(\F_q)\) to polynomiality in \(q\), positivity for split characters, monicity, and root properties of an associated quiver dimension vector. The supplied record leaves essential definitions and the scope of several assertions unresolved. Its partially solved status is supplied dated curation, not an independently verified conclusion.
\subsection{Sources}
\begin{itemize}
\item Local UnsolvedMath registry, supplied source record with permanent ID \(30001717\), exact source problem code \texttt{catalogue entry}, title \emph{Character Multiplicities and Kac--Moody Roots}, in the Oberwolfach Reports Open Problems collection. The record gives proposed year \(2011\) and names no individual proposer.
\item \emph{Representation Theory of Quivers and Finite Dimensional Algebras} (2011), Oberwolfach Reports, pp.\ 523--608, DOI \texttt{10.4171/owr/2011/10}: \url{https://doi.org/10.4171/owr/2011/10}. This is the report provenance supplied by the registry; no precise conjecture page or named contributor is supplied.
\item Supplied literature reference: Hausel--Letellier--Rodriguez-Villegas, \url{https://arxiv.org/abs/0810.2076}. The record supplies no theorem number identifying the results relevant to each assertion.
\item Supplied literature reference: Oberwolfach Report 10/2011, \url{https://doi.org/10.4171/owr/2011/10}.
\item Literature assessment checked at \texttt{2026-08-21}, preserved as supplied dated curation and not independently verified: the record labels the problem partially solved and states that the character-multiplicity assertions are proved for split semisimple characters and in indivisible cases through Kac-polynomial results, while the full general formulation in the extracted statement is not established. The supplied background expressly requests independent verification of time-sensitive status and cited claims.
\end{itemize}
\end{document}
